\section{Coefficient certificates and additional complete degrees}\label{sec:tables}
\subsection{All nonzero associative coefficients in degrees five and six}
Table~\ref{tab:certificate56} is a finite certificate for \cref{eq:Z5,eq:Z6}.
Its two coefficient columns on each side come from the cut formula
\cref{eq:word-cut} and expansion of the displayed right-nested brackets,
respectively. Both calculations give zero on every word absent from the table.
For degree five these are only $XXXXX$ and $YYYYY$; for degree six all $64$ possible
words, including the $36$ zero coefficients, are compared by the accompanying
script. Thus equality is checked in the free associative algebra, not in a
matrix quotient that might hide a Lie identity.

\begingroup\small
\renewcommand{\arraystretch}{1.45}
\begin{longtable}{@{}lrr@{\hspace{2.2em}}lrr@{}}
\caption{Complete nonzero word-coefficient certificates for the displayed fifth and sixth degrees.}\label{tab:certificate56}\\
\toprule
\multicolumn{3}{c}{$Z_5$} & \multicolumn{3}{c}{$Z_6$} \\

$w$ & $c_{\rm cut}$ & $c_{\rm bracket}$ & $w$ & $c_{\rm cut}$ & $c_{\rm bracket}$ \\
\midrule\endfirsthead
\multicolumn{6}{c}{\tablename\ \thetable\ (continued)}\\
\toprule
\multicolumn{3}{c}{$Z_5$} & \multicolumn{3}{c}{$Z_6$} \\

$w$ & $c_{\rm cut}$ & $c_{\rm bracket}$ & $w$ & $c_{\rm cut}$ & $c_{\rm bracket}$ \\
\midrule\endhead
\bottomrule\endfoot
$\mathrm{XXXXY}$ & $-\frac{1}{720}$ & $-\frac{1}{720}$ & $\mathrm{XXXXYY}$ & $-\frac{1}{1440}$ & $-\frac{1}{1440}$ \\
$\mathrm{XXXYX}$ & $\frac{1}{180}$ & $\frac{1}{180}$ & $\mathrm{XXXYXY}$ & $\frac{1}{360}$ & $\frac{1}{360}$ \\
$\mathrm{XXXYY}$ & $\frac{1}{180}$ & $\frac{1}{180}$ & $\mathrm{XXXYYY}$ & $\frac{1}{360}$ & $\frac{1}{360}$ \\
$\mathrm{XXYXX}$ & $-\frac{1}{120}$ & $-\frac{1}{120}$ & $\mathrm{XXYXXY}$ & $-\frac{1}{240}$ & $-\frac{1}{240}$ \\
$\mathrm{XXYXY}$ & $-\frac{1}{120}$ & $-\frac{1}{120}$ & $\mathrm{XXYXYY}$ & $-\frac{1}{240}$ & $-\frac{1}{240}$ \\
$\mathrm{XXYYX}$ & $-\frac{1}{120}$ & $-\frac{1}{120}$ & $\mathrm{XXYYXY}$ & $-\frac{1}{240}$ & $-\frac{1}{240}$ \\
$\mathrm{XXYYY}$ & $\frac{1}{180}$ & $\frac{1}{180}$ & $\mathrm{XXYYYY}$ & $-\frac{1}{1440}$ & $-\frac{1}{1440}$ \\
$\mathrm{XYXXX}$ & $\frac{1}{180}$ & $\frac{1}{180}$ & $\mathrm{XYXXXY}$ & $\frac{1}{360}$ & $\frac{1}{360}$ \\
$\mathrm{XYXXY}$ & $-\frac{1}{120}$ & $-\frac{1}{120}$ & $\mathrm{XYXXYY}$ & $-\frac{1}{240}$ & $-\frac{1}{240}$ \\
$\mathrm{XYXYX}$ & $\frac{1}{30}$ & $\frac{1}{30}$ & $\mathrm{XYXYXY}$ & $\frac{1}{60}$ & $\frac{1}{60}$ \\
$\mathrm{XYXYY}$ & $-\frac{1}{120}$ & $-\frac{1}{120}$ & $\mathrm{XYXYYY}$ & $\frac{1}{360}$ & $\frac{1}{360}$ \\
$\mathrm{XYYXX}$ & $-\frac{1}{120}$ & $-\frac{1}{120}$ & $\mathrm{XYYXXY}$ & $-\frac{1}{240}$ & $-\frac{1}{240}$ \\
$\mathrm{XYYXY}$ & $-\frac{1}{120}$ & $-\frac{1}{120}$ & $\mathrm{XYYXYY}$ & $-\frac{1}{240}$ & $-\frac{1}{240}$ \\
$\mathrm{XYYYX}$ & $\frac{1}{180}$ & $\frac{1}{180}$ & $\mathrm{XYYYXY}$ & $\frac{1}{360}$ & $\frac{1}{360}$ \\
$\mathrm{XYYYY}$ & $-\frac{1}{720}$ & $-\frac{1}{720}$ & $\mathrm{YXXXYX}$ & $-\frac{1}{360}$ & $-\frac{1}{360}$ \\
$\mathrm{YXXXX}$ & $-\frac{1}{720}$ & $-\frac{1}{720}$ & $\mathrm{YXXYXX}$ & $\frac{1}{240}$ & $\frac{1}{240}$ \\
$\mathrm{YXXXY}$ & $\frac{1}{180}$ & $\frac{1}{180}$ & $\mathrm{YXXYYX}$ & $\frac{1}{240}$ & $\frac{1}{240}$ \\
$\mathrm{YXXYX}$ & $-\frac{1}{120}$ & $-\frac{1}{120}$ & $\mathrm{YXYXXX}$ & $-\frac{1}{360}$ & $-\frac{1}{360}$ \\
$\mathrm{YXXYY}$ & $-\frac{1}{120}$ & $-\frac{1}{120}$ & $\mathrm{YXYXYX}$ & $-\frac{1}{60}$ & $-\frac{1}{60}$ \\
$\mathrm{YXYXX}$ & $-\frac{1}{120}$ & $-\frac{1}{120}$ & $\mathrm{YXYYXX}$ & $\frac{1}{240}$ & $\frac{1}{240}$ \\
$\mathrm{YXYXY}$ & $\frac{1}{30}$ & $\frac{1}{30}$ & $\mathrm{YXYYYX}$ & $-\frac{1}{360}$ & $-\frac{1}{360}$ \\
$\mathrm{YXYYX}$ & $-\frac{1}{120}$ & $-\frac{1}{120}$ & $\mathrm{YYXXXX}$ & $\frac{1}{1440}$ & $\frac{1}{1440}$ \\
$\mathrm{YXYYY}$ & $\frac{1}{180}$ & $\frac{1}{180}$ & $\mathrm{YYXXYX}$ & $\frac{1}{240}$ & $\frac{1}{240}$ \\
$\mathrm{YYXXX}$ & $\frac{1}{180}$ & $\frac{1}{180}$ & $\mathrm{YYXYXX}$ & $\frac{1}{240}$ & $\frac{1}{240}$ \\
$\mathrm{YYXXY}$ & $-\frac{1}{120}$ & $-\frac{1}{120}$ & $\mathrm{YYXYYX}$ & $\frac{1}{240}$ & $\frac{1}{240}$ \\
$\mathrm{YYXYX}$ & $-\frac{1}{120}$ & $-\frac{1}{120}$ & $\mathrm{YYYXXX}$ & $-\frac{1}{360}$ & $-\frac{1}{360}$ \\
$\mathrm{YYXYY}$ & $-\frac{1}{120}$ & $-\frac{1}{120}$ & $\mathrm{YYYXYX}$ & $-\frac{1}{360}$ & $-\frac{1}{360}$ \\
$\mathrm{YYYXX}$ & $\frac{1}{180}$ & $\frac{1}{180}$ & $\mathrm{YYYYXX}$ & $\frac{1}{1440}$ & $\frac{1}{1440}$ \\
$\mathrm{YYYXY}$ & $\frac{1}{180}$ & $\frac{1}{180}$ & -- & -- & -- \\
$\mathrm{YYYYX}$ & $-\frac{1}{720}$ & $-\frac{1}{720}$ & -- & -- & -- \\
\end{longtable}
\endgroup

\subsection{Complete seventh and eighth BCH degrees}
In the following table the entries specify exactly
$Z_n=\sum_w c_wL(w)$ for $n=7,8$, using \cref{eq:lyndon-definition};
there are no omitted nonzero terms in this convention. Expanding the brackets
and applying \cref{eq:word-cut} proves each identity by the same finite procedure
as in the preceding table. The CSV file continues these full expansions through
degree twelve.

\begingroup\small
\renewcommand{\arraystretch}{1.45}
\begin{longtable}{@{}lr@{\hspace{3em}}lr@{}}
\caption{Additional complete BCH homogeneous polynomials in standard Lyndon bracketing.}\label{tab:bch78}\\
\toprule
\multicolumn{2}{c}{$Z_7$} & \multicolumn{2}{c}{$Z_8$} \\

$w$ & coefficient of $L(w)$ & $w$ & coefficient of $L(w)$ \\
\midrule\endfirsthead
\multicolumn{4}{c}{\tablename\ \thetable\ (continued)}\\
\toprule
\multicolumn{2}{c}{$Z_7$} & \multicolumn{2}{c}{$Z_8$} \\

$w$ & coefficient of $L(w)$ & $w$ & coefficient of $L(w)$ \\
\midrule\endhead
\bottomrule\endfoot
$\mathrm{XXXXXXY}$ & $\frac{1}{30240}$ & $\mathrm{XXXXXXYY}$ & $\frac{1}{60480}$ \\
$\mathrm{XXXXXYY}$ & $-\frac{1}{5040}$ & $\mathrm{XXXXXYXY}$ & $-\frac{1}{15120}$ \\
$\mathrm{XXXXYXY}$ & $\frac{1}{10080}$ & $\mathrm{XXXXXYYY}$ & $-\frac{1}{10080}$ \\
$\mathrm{XXXXYYY}$ & $\frac{1}{3780}$ & $\mathrm{XXXXYXXY}$ & $\frac{1}{20160}$ \\
$\mathrm{XXXYXXY}$ & $\frac{1}{10080}$ & $\mathrm{XXXXYXYY}$ & $-\frac{1}{20160}$ \\
$\mathrm{XXXYXYY}$ & $\frac{1}{1680}$ & $\mathrm{XXXXYYXY}$ & $\frac{1}{2520}$ \\
$\mathrm{XXXYYXY}$ & $\frac{1}{1260}$ & $\mathrm{XXXXYYYY}$ & $\frac{23}{120960}$ \\
$\mathrm{XXXYYYY}$ & $\frac{1}{3780}$ & $\mathrm{XXXYXXYY}$ & $\frac{1}{4032}$ \\
$\mathrm{XXYXXYY}$ & $\frac{1}{2016}$ & $\mathrm{XXXYXYXY}$ & $-\frac{1}{10080}$ \\
$\mathrm{XXYXYXY}$ & $-\frac{1}{5040}$ & $\mathrm{XXXYXYYY}$ & $\frac{13}{30240}$ \\
$\mathrm{XXYXYYY}$ & $\frac{13}{15120}$ & $\mathrm{XXXYYXYY}$ & $\frac{1}{20160}$ \\
$\mathrm{XXYYXYY}$ & $\frac{1}{10080}$ & $\mathrm{XXXYYYXY}$ & $-\frac{1}{3024}$ \\
$\mathrm{XXYYYXY}$ & $-\frac{1}{1512}$ & $\mathrm{XXXYYYYY}$ & $-\frac{1}{10080}$ \\
$\mathrm{XXYYYYY}$ & $-\frac{1}{5040}$ & $\mathrm{XXYXYXYY}$ & $\frac{1}{2520}$ \\
$\mathrm{XYXYXYY}$ & $\frac{1}{1260}$ & $\mathrm{XXYXYYYY}$ & $-\frac{1}{4032}$ \\
$\mathrm{XYXYYYY}$ & $-\frac{1}{2016}$ & $\mathrm{XXYYXYYY}$ & $-\frac{1}{10080}$ \\
$\mathrm{XYYXYYY}$ & $-\frac{1}{5040}$ & $\mathrm{XXYYYYYY}$ & $\frac{1}{60480}$ \\
$\mathrm{XYYYYYY}$ & $\frac{1}{30240}$ & -- & -- \\
\end{longtable}
\endgroup

\subsection{Complete fifth and sixth Zassenhaus exponents}
These are the next two full exponents after the terms displayed on the source
page. A coefficient multiplies $L(w)$, not $R(w)$. Every entry follows either
from \cref{eq:zass-residual} or from \cref{eq:zass-finite-Lie}; the two constructions
agree exactly. Higher exponents through degree twelve are supplied as CSV data.

\begingroup\small
\renewcommand{\arraystretch}{1.45}
\begin{longtable}{@{}lr@{\hspace{3em}}lr@{}}
\caption{Complete fifth and sixth Zassenhaus exponents.}\label{tab:zass56}\\
\toprule
\multicolumn{2}{c}{$C_5$} & \multicolumn{2}{c}{$C_6$} \\

$w$ & coefficient of $L(w)$ & $w$ & coefficient of $L(w)$ \\
\midrule\endfirsthead
\multicolumn{4}{c}{\tablename\ \thetable\ (continued)}\\
\toprule
\multicolumn{2}{c}{$C_5$} & \multicolumn{2}{c}{$C_6$} \\

$w$ & coefficient of $L(w)$ & $w$ & coefficient of $L(w)$ \\
\midrule\endhead
\bottomrule\endfoot
$\mathrm{XXXXY}$ & $\frac{1}{120}$ & $\mathrm{XXXXXY}$ & $-\frac{1}{720}$ \\
$\mathrm{XXXYY}$ & $-\frac{1}{30}$ & $\mathrm{XXXXYY}$ & $\frac{1}{144}$ \\
$\mathrm{XXYXY}$ & $-\frac{1}{60}$ & $\mathrm{XXXYYY}$ & $-\frac{1}{72}$ \\
$\mathrm{XXYYY}$ & $\frac{1}{20}$ & $\mathrm{XXYYYY}$ & $\frac{1}{72}$ \\
$\mathrm{XYXYY}$ & $-\frac{1}{20}$ & $\mathrm{XYXYYY}$ & $-\frac{1}{72}$ \\
$\mathrm{XYYYY}$ & $-\frac{1}{30}$ & $\mathrm{XYYYYY}$ & $-\frac{1}{144}$ \\
\end{longtable}
\endgroup

\subsection{Sizes of the exported expansions}
\begin{table}[htbp]
\centering\small
\begin{tabular}{@{}rrrr@{}}
\toprule
Degree & BCH words & BCH Lyndon entries & Zassenhaus Lyndon entries\\
\midrule

1 & 2 & 2 & 0 \\
2 & 2 & 1 & 1 \\
3 & 6 & 2 & 2 \\
4 & 4 & 1 & 3 \\
5 & 30 & 6 & 6 \\
6 & 28 & 5 & 6 \\
7 & 126 & 18 & 18 \\
8 & 124 & 17 & 29 \\
9 & 390 & 55 & 56 \\
10 & 388 & 55 & 95 \\
11 & 2046 & 186 & 186 \\
12 & 2044 & 185 & 326 \\
\bottomrule
\end{tabular}
\caption{Numbers of nonzero coefficients in the supplied encodings. The degree-one Zassenhaus factor $e^Y$ is separated from the exponents $C_n$, $n\geq2$.}\label{tab:counts}
\end{table}
